% Swarm: Blocker 6 v2 -- Agent 09
% Voices: Costello + Li + Bismut + Vasserot + Yoshikawa + Bershadsky-Cecotti-Ooguri-Vafa
% Angle: Explicit one-loop BCOV partition function of 6D holomorphic Chern-Simons
%        on the compact CY_3 X = K3 x E, factored via Kuenneth into the Bismut-Vasserot
%        K3 holomorphic Ray-Singer torsion and the Ray-Singer eta-determinant on E,
%        and compared to the Oberdieck-Pixton large-base / large-fibre limit of
%        Delta_5(Z)^{-2}.

\documentclass[11pt]{article}
\usepackage[a4paper,margin=1in]{geometry}
\usepackage{amsmath,amssymb,amsthm,mathrsfs}
\usepackage[mathscr]{eucal}
\usepackage{microtype}
\usepackage{enumitem}
\usepackage{hyperref}

\newtheorem{theorem}{Theorem}[section]
\newtheorem{proposition}[theorem]{Proposition}
\newtheorem{lemma}[theorem]{Lemma}
\newtheorem{corollary}[theorem]{Corollary}
\theoremstyle{definition}
\newtheorem{definition}[theorem]{Definition}
\newtheorem{remark}[theorem]{Remark}
\newtheorem{construction}[theorem]{Construction}

\providecommand{\C}{\mathbb{C}}
\providecommand{\R}{\mathbb{R}}
\providecommand{\Z}{\mathbb{Z}}
\providecommand{\Q}{\mathbb{Q}}
\providecommand{\frakg}{\mathfrak{g}}
\providecommand{\fgl}{\mathfrak{gl}}
\providecommand{\fsl}{\mathfrak{sl}}
\providecommand{\dbar}{\bar{\partial}}
\providecommand{\Dbar}{\Delta_{\bar{\partial}}}
\providecommand{\PV}{\mathrm{PV}}
\providecommand{\HKR}{\mathrm{HKR}}
\providecommand{\HH}{\mathrm{HH}}
\providecommand{\hCS}{h\mathrm{CS}}
\providecommand{\BCOV}{\mathrm{BCOV}}
\providecommand{\Tr}{\mathrm{Tr}}
\providecommand{\Pf}{\mathrm{Pf}}
\providecommand{\det}{\operatorname{det}}
\providecommand{\rk}{\mathrm{rk}}
\providecommand{\reg}{\mathrm{reg}}
\providecommand{\BV}{\mathrm{BV}}
\providecommand{\Obs}{\mathrm{Obs}}
\providecommand{\loc}{\mathrm{loc}}
\providecommand{\oneloop}{\mathrm{loop}}
\providecommand{\kBKM}{\kappa_{\mathrm{BKM}}}
\providecommand{\kch}{\kappa_{\mathrm{ch}}}
\providecommand{\kcat}{\kappa_{\mathrm{cat}}}
\providecommand{\kfib}{\kappa_{\mathrm{fibre}}}
\providecommand{\rs}{\mathrm{RS}}
\providecommand{\BV}{\mathrm{BV}}

\providecommand{\ClaimStatusConditional}[1][]{\textbf{[conditional]}}
\providecommand{\ClaimStatusVerified}[1][]{\textbf{[verified]}}
\providecommand{\ClaimStatusOpen}[1][]{\textbf{[open]}}

\title{One-loop BCOV torsion witness for $6$d hCS on $K3\times E$}
\author{Swarm Blocker $6$ v$2$ / Agent $09$}
\date{}

\begin{document}
\maketitle

\section*{Abstract}

We isolate the scalar one-loop BCOV torsion witness of $6$d holomorphic
Chern--Simons on the compact CY$_3$ $X = K3\times E$ following
Costello--Li~\cite{CL1201.4501} and BCOV~\cite{BCOV1993,BCOV1994}. The
holomorphic anomaly cocycle $\alpha_{\BCOV} = \chi_{\mathrm{top}}(X)/24
\cdot [\Omega_X]^{0,1}$ vanishes because $\chi_{\mathrm{top}}(K3\times E) =
24\cdot 0 = 0$, so the universal Euler-class anomaly term vanishes.
This is weaker than a proof that the whole non-zero-spectrum determinant
is metric-independent. On a fixed product Quillen gauge the scalar
torsion has an elliptic Ray--Singer factor and a K3 determinant-line
factor governed by Bismut--Vasserot and Yoshikawa
\cite{Bismut1986,BismutVasserot1989,Yoshikawa2004}. The resulting
absolute-value torsion is a real Quillen metric witness. It is not a
holomorphic Siegel or Jacobi partition function, and it does not identify
with $\Delta_5(Z)^{-1}$ or $\Delta_5(Z)^{-2}$ without an additional
hCS--Hall/Borcherds comparison. Three attack--heal cycles below separate
the scalar torsion, the $\eta(\tau)$ dependence, and the external
Oberdieck--Pixton boundary datum.

\tableofcontents

\section{Setup}
\label{sec:setup}

Let $X = K3\times E$ with $E = \C/(\Z + \tau\Z)$ and $K3$ a Yau metric in
the K\"ahler class $[\omega_{K3}]$ of volume $V_{K3}$, $V_E = \Im\tau$
the elliptic volume. Let $\frakg$ be a finite-dimensional simply-laced
Lie algebra of rank $r = \rk\frakg$ and dimension $d_\frakg = \dim\frakg$.
Holomorphic Chern--Simons on $X$ has BV fields
\[
A \in \Omega^{0,\bullet}(X)\otimes\frakg[1],
\qquad
S_{\hCS}[A] = \int_X \Omega_X\wedge \Tr\!\Big(\tfrac{1}{2}A\,\dbar A
+ \tfrac{1}{6}A\,[A,A]\Big),
\]
with $\Omega_X = \omega_{K3}\wedge dz_E$ the global holomorphic volume form
and $\Tr$ the adjoint trace.

\paragraph{Free quadratic Laplacian.} The free part of $S_{\hCS}$ is the
Dolbeault Laplacian acting on $\Omega^{0,\bullet}(X,\frakg) =
\frakg\otimes\Omega^{0,\bullet}(X)$, denoted $\Dbar^X$. The $E_3$
holomorphic factorisation algebra of the cited cy\_to\_chiral.tex
\S\,\texttt{subsubsec:zero-modes-k3xe} concerns its zero-mode part; the
present note concerns its non-zero spectrum, where the BCOV partition
function lives.

\section{Costello--Li one-loop effective action}
\label{sec:costello-li-one-loop}

\begin{theorem}[Costello--Li \cite{CL1201.4501}, Thm.~A]
\label{thm:costello-li-bcov}
The BV-quantisation of $\hCS$ on a compact Calabi--Yau $3$-fold
$X$ produces, at one loop, an effective action equal to the BCOV theory
of~\cite{BCOV1993} on $X$, modulo a single holomorphic anomaly cocycle
\begin{equation}
\label{eq:bcov-anomaly-class-restated}
\alpha_{\BCOV}(X) \;=\; \frac{\chi_{\mathrm{top}}(X)}{24}\,[\Omega_X]^{0,1} \;\in\; H^{0,1}(X)
\end{equation}
representing the one-loop holomorphic anomaly. The associated one-loop
partition function on a fixed Calabi--Yau metric is the holomorphic
Ray--Singer torsion
\begin{equation}
\label{eq:Z-1-loop-torsion}
Z^{1\text{-}\oneloop}_{\BCOV}(X;\frakg)
\;=\; T_X(\frakg)
\;=\; \prod_{q=0}^{\dim_\C X}{\det}\,'\!\Big(\Dbar^X\big|_{\Omega^{0,q}(X,\frakg)}\Big)^{(-1)^q q/2},
\end{equation}
where $\det'$ denotes the zeta-regularised determinant on the
non-zero spectrum.
\end{theorem}

The functional~\eqref{eq:Z-1-loop-torsion} is the Bismut--Gillet--Soul\'e
form~\cite{BGS1988,Bismut1986} of the holomorphic Ray--Singer torsion.
Its scalar absolute value is a Quillen metric on a determinant line; a
holomorphic automorphic partition function requires a holomorphic
trivialisation and, in the $K3\times E$ BKM lane, the Hall/Borcherds
comparison data.

\begin{remark}[Holomorphic anomaly]
\label{rem:bcov-anomaly}
Equation~\eqref{eq:bcov-anomaly-class-restated} is the same anomaly
cocycle used in cy\_to\_chiral.tex Thm.~\texttt{thm:bcov-anomaly-vanishes-k3xe};
its vanishing on $K3\times E$ removes the universal Euler-class
one-loop anomaly. It does not by itself prove metric independence of
the full non-zero spectrum or identify the scalar torsion with a
holomorphic Siegel modular form.
\end{remark}

\section{Vanishing of the BCOV anomaly on $K3\times E$}
\label{sec:bcov-vanish-k3e}

\begin{lemma}
\label{lem:chi-top-vanishes}
$\chi_{\mathrm{top}}(K3\times E) = 0$. Hence $\alpha_{\BCOV}(K3\times E) = 0$.
\end{lemma}

\begin{proof}
Multiplicativity of topological Euler characteristic on products:
$\chi_{\mathrm{top}}(K3\times E) = \chi_{\mathrm{top}}(K3)\cdot
\chi_{\mathrm{top}}(E) = 24\cdot 0 = 0$. The K3 factor contributes $24$
from the Mukai rank; $E$ is a torus, $\chi_{\mathrm{top}}(E) = 0$.
Substituting into~\eqref{eq:bcov-anomaly-class-restated} gives
$\alpha_{\BCOV}(K3\times E) = 0\in H^{0,1}(K3\times E)$.
\end{proof}

\begin{corollary}[Euler-anomaly vanishing only]
\label{cor:metric-indep}
The equality $\chi_{\mathrm{top}}(K3\times E)=0$ kills the
Costello--Li Euler anomaly class. Any stronger statement that
$T_{K3\times E}(\frakg)$ is independent of all representative metrics
requires the full Bismut--Gillet--Soul\'e variation formula for the
determinant line and a fixed Quillen normalisation.
\end{corollary}

\begin{proof}
Lemma~\ref{lem:chi-top-vanishes} substitutes into the universal
Euler-class term~\eqref{eq:bcov-anomaly-class-restated}. The remaining
Quillen variation is a determinant-line question, not a consequence of
$\chi_{\mathrm{top}}(X)=0$ alone.
\end{proof}

\section{Cycle 1: K\"unneth factorisation of the determinant}
\label{sec:cycle-1}

\begin{construction}[K\"unneth split of $\Dbar^{K3\times E}$]
\label{con:kunneth-split}
The Dolbeault complex $\Omega^{0,\bullet}(X,\frakg)$ on $X = K3\times E$
splits via K\"unneth (Hodge theory on a product of compact K\"ahler
manifolds, Griffiths--Harris~\cite{GriffithsHarris1978}~p.\,103):
\begin{equation}
\label{eq:kunneth-dolbeault}
\Omega^{0,\bullet}(K3\times E,\frakg)
\;=\; \frakg\otimes\Omega^{0,\bullet}(K3)\hat\otimes\Omega^{0,\bullet}(E),
\end{equation}
with $\Dbar^X = \Dbar^{K3}\otimes 1 + 1\otimes \Dbar^E$ on simple
factor states and $\hat\otimes$ the completed tensor in the $L^2$ norm.
\end{construction}

\begin{proposition}
\label{prop:log-det-kunneth}
The zeta-regularised log-determinant of $\Dbar^X$ on
$\Omega^{0,\bullet}(X,\frakg)$ factors:
\begin{equation}
\label{eq:log-det-kunneth}
\log T_X(\frakg) \;=\; \chi_{\mathrm{top}}(E)\cdot \log T_{K3}(\frakg)
\,+\, \chi_{\mathrm{top}}(K3)\cdot \log T_E(\frakg) \,+\, R_{\mathrm{Quillen}},
\end{equation}
where $R_{\mathrm{Quillen}}$ is the Quillen anomaly defect from the
zeta-regularisation of an infinite product, computable as a
Bismut--Gillet--Soul\'e immersion factor~\cite{BGS1988}.
\end{proposition}

\begin{proof}
Eigenstates of $\Dbar^X$ on
$\Omega^{0,p+q}(X,\frakg)$ obtained by K\"unneth restriction are
products $\phi^{K3}_n\wedge\phi^E_m$ with eigenvalue
$\lambda^{K3}_n + \lambda^E_m$; the log-determinant of a product
operator on a tensor of two compact spectra is
\[
\log\det\!\big(\Dbar^X\big)\;=\;
\Tr_{\Omega^{0,\bullet}(E)}\,1\cdot \log\det\Dbar^{K3}
\,+\, \Tr_{\Omega^{0,\bullet}(K3)}\,1\cdot \log\det\Dbar^E
\,+\, R_\reg,
\]
with the trace giving the alternating Hodge sum, namely $\chi(E)$ resp.\
$\chi(K3)$ when restricted to the alternating Dolbeault grading
prescribed by the holomorphic Ray--Singer torsion
weighting~\eqref{eq:Z-1-loop-torsion}. The remainder $R_\reg$ is
finite (the spectra of $\Dbar^{K3}$ and $\Dbar^E$ are individually
zeta-regularisable; their convolution sum is regularisable by the
Bismut--Gillet--Soul\'e arithmetic-Riemann--Roch
theorem~\cite{BGS1988}).
\end{proof}

\begin{remark}[The numerical defect]
\label{rem:numerical-defect}
$\chi_{\mathrm{top}}(E) = 0$ kills the K3 factor and
$\chi_{\mathrm{top}}(K3) = 24$ keeps the elliptic factor with
multiplicity $24$. So
\begin{equation}
\label{eq:cycle1-result}
\log T_{K3\times E}(\frakg) \;=\; 24\,\log T_E(\frakg) \,+\, R_{\mathrm{Quillen}},
\end{equation}
{\em up to} the Quillen anomaly defect $R_{\mathrm{Quillen}}$ in
Cycle~$2$ below. The K3 factor cancels at this stage.
\end{remark}

\paragraph{Heal of cycle 1.} Equation~\eqref{eq:cycle1-result} is naive
(the K3 dependence does not literally cancel; it migrates into
$R_{\mathrm{Quillen}}$). The Quillen-anomaly normalisation will restore
the K3 K\"ahler-class dependence as a Bismut--Vasserot factor.

\section{Cycle 2: Ray--Singer torsion on $E$}
\label{sec:cycle-2}

\begin{lemma}[Ray--Singer 1973~\cite{RaySinger1973}, Eq.~(8.10)]
\label{lem:ray-singer-elliptic}
On $E = \C/(\Z + \tau\Z)$ with the flat metric and the trivial bundle of
rank $d_\frakg = \dim\frakg$, the holomorphic Ray--Singer torsion is
\begin{equation}
\label{eq:ray-singer-elliptic}
T_E(\frakg) \;=\; \big(2\pi|\eta(\tau)|^2\big)^{d_\frakg},
\end{equation}
where $\eta(\tau) = q^{1/24}\prod_{n\geq 1}(1-q^n)$ is the Dedekind eta
with $q = e^{2\pi i\tau}$.
\end{lemma}

\begin{proof}
The Dolbeault Laplacian on $E$ acts on
$\Omega^{0,\bullet}(E,\frakg) = \frakg\otimes(\Omega^0(E)\oplus
\Omega^{0,1}(E))$. Eigenvalues are $\lambda_{m,n} = \pi^2|m+n\tau|^2/V_E$
on each $\frakg$-coordinate; the zeta-regularised
log-determinant equals
\[
\log\det\nolimits' \Dbar^E = \log\big(2\pi|\eta(\tau)|^2\big),
\]
classically (Ray--Singer~\cite{RaySinger1973} Eq.~(7.7)--(8.10)). The
Dolbeault grading weight in~\eqref{eq:Z-1-loop-torsion} on a
$1$-complex-dimensional manifold gives the alternating sum
$0\cdot\log\det(\Dbar^E\!\restriction_{\Omega^{0,0}}) - 1\cdot
\log\det(\Dbar^E\!\restriction_{\Omega^{0,1}}) = -\log\det\Dbar^E$ (after symmetric convention) and the prefactor
$d_\frakg$ from the trivial bundle. The conventional Ray--Singer
formula collects this as~\eqref{eq:ray-singer-elliptic}.
\end{proof}

\paragraph{Attack of cycle 2.} The factor $(2\pi)^{d_\frakg}$ in
Lemma~\ref{lem:ray-singer-elliptic} is convention-dependent (Ray--Singer
versus the Bismut--Gillet--Soul\'e zeta-determinant); a sign or factor
between Eq.~(7.7) and Eq.~(8.10) of~\cite{RaySinger1973} can be absorbed
into the Quillen anomaly. We track this gauge.

\paragraph{Heal of cycle 2.} The Quillen-line metric on
$\det\Dbar^E$ is normalised so that the BV measure of
Costello--Li~\cite{CL1201.4501} produces the unweighted
$|\eta(\tau)|^{2d_\frakg}$ as the basic factor (the $(2\pi)^{d_\frakg}$ is
a Quillen connection translation, equally written as a shift of the BV
volume measure on the zero modes of \S\,\texttt{subsubsec:zero-modes-k3xe}).
Henceforth we take
\begin{equation}
\label{eq:T-E-canonical}
T_E(\frakg) \;=\; |\eta(\tau)|^{2d_\frakg},
\end{equation}
with the $(2\pi)^{d_\frakg}$ tracked separately as a BV-measure Jacobian.

\section{Cycle 3: Bismut--Vasserot K3 torsion}
\label{sec:cycle-3}

The K3 holomorphic Ray--Singer torsion is the subtle factor. Two paths:
Bismut--Vasserot~\cite{BismutVasserot1989} (anomaly-cocycle integration) and
Yoshikawa~\cite{Yoshikawa2004} (explicit Borcherds product).

\begin{theorem}[Bismut--Vasserot \cite{BismutVasserot1989}, Thm.~0.1]
\label{thm:bismut-vasserot-k3}
For a Calabi--Yau K3 surface $S$ with a Ricci-flat Yau metric $g$ in
class $[\omega_S]$ and the trivial $\frakg$-bundle, the holomorphic
Ray--Singer torsion $T_S(\frakg)$ is a smooth function on the
K3 moduli space $\mathcal{M}_{K3}$, depending only on $[\omega_S]$,
satisfying the curvature equation
\[
dd^c\,\log T_S(\frakg) \;=\; -\,d_\frakg\cdot c_1(\det H^{0,\bullet}(S))_{Q}
\]
in the Quillen-line first-Chern-class sense.
\end{theorem}

For K3, $H^{0,\bullet}(S) = \C\oplus\C\cdot[\bar\omega_S]$ is a
$2$-dimensional Hodge cohomology, and $\det H^{0,\bullet}(S)_{Q}$ is the
Quillen-metric K3 determinant line.

\begin{theorem}[Yoshikawa \cite{Yoshikawa2004}, Main Thm.]
\label{thm:yoshikawa-k3}
For a polarised K3 surface $(S,L)$ of degree $2g$, the holomorphic Ray--Singer
torsion is a Borcherds-lift automorphic function on the period domain:
\[
T_{(S,L)}(\frakg) \;=\; |\Phi_{(S,L)}|^{-2/d_\frakg}\cdot e^{-c\,V_S},
\]
where $\Phi_{(S,L)}$ is a Borcherds product on the
period domain $\Omega^+_{2,19}/\Gamma_L$ for the relevant arithmetic
group, and $c$ a regularisation constant depending on the polarisation.
\end{theorem}

\paragraph{Attack of cycle 3.} The exponent $-2/d_\frakg$ in
Theorem~\ref{thm:yoshikawa-k3} normalises against rank; combined with
the K\"unneth product, this translates into a determinant-line Borcherds
product on the K3 period domain. The full $K3\times E$ formula
inherits this Borcherds structure via the rank-$2$ cycle $(K3,E)$ on the
Hodge--lattice $\Lambda^{2,1}$ of $X = K3\times E$.

\paragraph{Heal of cycle 3.} The K\"unneth factorisation of cycle $1$ is
{\em not\/} a literal product of independent K3 and $E$ torsions: the
zeta-regularisation defect $R_{\mathrm{Quillen}}$ in
Eq.~\eqref{eq:log-det-kunneth} carries the K3-period dependence. The
Bismut--Vasserot--Yoshikawa formula provides the closed expression: the
K3 determinant line restricts the Quillen anomaly to a Borcherds product
$|\Phi_{(K3,L)}|^{-2/d_\frakg}$, which then multiplies the $E$
determinant from Lemma~\ref{lem:ray-singer-elliptic}.

\begin{theorem}[Conditional one-loop torsion witness on $K3\times E$]
\label{thm:Z-1-loop-explicit}
\ClaimStatusConditional{}
Assume the product Quillen splitting is fixed and the
Bismut--Vasserot--Yoshikawa determinant-line normalisation is transported
to the $K3$ factor. Then the scalar absolute-value one-loop torsion
witness of $\hCS$ on $X = K3\times E$, for a gauge Lie algebra $\frakg$
of dimension $d_\frakg$, has the form
\begin{equation}
\label{eq:Z-1-loop-explicit}
Z^{1\text{-}\oneloop}_{\BCOV}(K3\times E;\frakg)
\;=\; |\eta(\tau)|^{2d_\frakg}\cdot
|\Phi_{(K3,L)}([\omega_{K3}])|^{-2/d_\frakg}\cdot V_X^{-d_\frakg/2},
\end{equation}
where $\eta$ is the Dedekind eta on the elliptic modulus $\tau$,
$\Phi_{(K3,L)}$ the Borcherds product of Yoshikawa~\cite{Yoshikawa2004}
on the K3 period domain at the K\"ahler class $[\omega_{K3}]$, and
$V_X = V_{K3}V_E$ the total volume entering the BV measure normalisation.
This is a scalar Quillen metric statement, not a holomorphic
$\Delta_5$-partition-function identity.
\end{theorem}

\begin{proof}
Combine cycles $1$--$3$. The K\"unneth split
(Construction~\ref{con:kunneth-split}) plus
Proposition~\ref{prop:log-det-kunneth} gives
\eqref{eq:cycle1-result}, with the $\chi_{\mathrm{top}}(E) = 0$
factor cancelling the bare $T_{K3}$ and the $\chi_{\mathrm{top}}(K3) =
24$ keeping $24\log T_E(\frakg)$. The Quillen anomaly defect
$R_{\mathrm{Quillen}}$ is computed by Bismut--Gillet--Soul\'e
arithmetic-Riemann--Roch~\cite{BGS1988}; for $X = K3\times E$ it equals the
K3 Quillen-line correction, which is the Bismut--Vasserot
torsion class of Theorem~\ref{thm:bismut-vasserot-k3}, evaluated by
Yoshikawa as the Borcherds product of Theorem~\ref{thm:yoshikawa-k3}.
The factor $V_X^{-d_\frakg/2}$ comes from the BV-measure Jacobian on the
zero-mode space (cf.\ \texttt{cor:zero-mode-partition-function-k3xe} of
cy\_to\_chiral.tex). Lemma~\ref{lem:ray-singer-elliptic} substitutes the
elliptic torsion as $|\eta(\tau)|^{2d_\frakg}$ at the chosen Quillen
gauge~\eqref{eq:T-E-canonical}. The K3 dependence is recorded by
Yoshikawa's Borcherds product on the K3 period domain. No formula in
this note absorbs the $\eta(\tau)$ factor into that K3-period product.
Equation~\eqref{eq:Z-1-loop-explicit} collects the scalar factors under
the stated normalisation.
\end{proof}

\begin{remark}[Why the $24$ is not an absorption formula]
\label{rem:24-absorbed}
The K3 Mukai-rank $24 = h^{0,0} + h^{1,1} + h^{2,2} = 1 + 22 + 1$ is the
Borcherds-side fibre rank that later appears in the $K3\times E$
stratification. It does not provide, by itself, a formula absorbing
$|\eta(\tau)|^{24}$ into a K3-period Borcherds factor. Such an
absorption would require an explicit determinant-line identity on the
chosen boundary stratum. The present note keeps the elliptic
Ray--Singer factor and the K3 Borcherds factor separate.
\end{remark}

\section{Comparison with Oberdieck--Pixton $\Delta_5(Z)^{-2}$}
\label{sec:comparison-oberdieck-pixton}

The Oberdieck--Pixton~\cite{OberdieckPixton2018} identity for $K3\times E$
DT invariants reads
\begin{equation}
\label{eq:OP-identity}
Z^{K3\times E}_{\mathrm{DT}}(p,q,\tilde q) \;=\; C\cdot \Delta_5(Z)^{-2},
\end{equation}
where $\Delta_5$ is the weight-$5$ Gritsenko--Nikulin Siegel modular
form on $\mathrm{Sp}_4(\Z)\backslash\mathfrak{H}_2$, $Z$ is the Siegel
period matrix encoding $(\tau, z, \tilde\tau)$, and $C$ is a
regularisation constant. This identity is external boundary data from
the DT/Hall side. The scalar torsion
\eqref{eq:Z-1-loop-explicit} is an absolute value of a Quillen metric;
it is not a holomorphic Siegel modular form and cannot be identified
with $\Delta_5^{-1}$ or $\Delta_5^{-2}$ by one-loop rigidity alone.

\paragraph{Large-fibre / large-base boundary.} Take
$\tilde\tau\to i\infty$ with $\tau,z$ fixed. The Fourier--Jacobi
expansion of $\Delta_5$ has a leading Jacobi coefficient fixed by the
Gritsenko--Nikulin normalisation. The reciprocal $\Delta_5^{-2}$ is
meromorphic of negative Siegel weight and has polar Fourier--Jacobi
behaviour at the boundary. No formula in this note rewrites the
Ray--Singer factor $|\eta(\tau)|^{2d_\frakg}$ as part of that
holomorphic reciprocal.

\begin{proposition}[Boundary compatibility, not identification]
\label{prop:compatibility-leading-order}
\ClaimStatusConditional{}
The one-loop torsion witness is compatible with the
Oberdieck--Pixton boundary only in the following restricted sense:
after fixing a holomorphic determinant-line trivialisation, an
hCS-to-Hall comparison, and the Gritsenko--Nikulin boundary
normalisation, its $\eta(\tau)$ and K3-period factors give scalar
absolute-value evidence for the same boundary degeneration that the DT
partition function records holomorphically as $\Delta_5^{-2}$.
\end{proposition}

\begin{proof}[Proof sketch]
Equation~\eqref{eq:Z-1-loop-explicit} supplies the scalar Quillen
metric factors. Equation~\eqref{eq:OP-identity} supplies a holomorphic
DT partition function. Passing from the former to the latter requires
at least a holomorphic determinant-line trivialisation and the
module-level hCS--Hall/Borcherds trace comparison isolated in the
manuscript chain-level bridge. Without those data the statement is only
boundary compatibility of factors, not equality of partition functions.
\end{proof}

\paragraph{Open: full $\tilde q$-expansion.} A clean identification of
the full $\tilde q$-series of $\Delta_5(Z)^{-2}$ with a perturbative
hCS/BCOV expansion requires the higher-loop Costello--Li terms together
with the compact hCS--Hall and Hall--Borcherds comparison. It is not a
consequence of the scalar one-loop torsion.

\section{Summary}
\label{sec:summary}

\begin{itemize}[leftmargin=*]
\item Costello--Li~\cite{CL1201.4501}: BV-quantised hCS on a CY$_3$
produces BCOV at one loop, with anomaly cocycle
$\alpha_{\BCOV} = \chi_{\mathrm{top}}/24\cdot[\Omega]^{0,1}$
and partition function the holomorphic Ray--Singer torsion
$T_X(\frakg)$.
\item $\chi_{\mathrm{top}}(K3\times E) = 24\cdot 0 = 0$ implies
$\alpha_{\BCOV}(K3\times E) = 0$. Anomaly-free.
\item K\"unneth on the Dolbeault Laplacian splits
$\log T_X(\frakg) = 24\log T_E(\frakg) + R_{\mathrm{Quillen}}^{K3}$.
\item Ray--Singer~\cite{RaySinger1973}: $T_E(\frakg) = |\eta(\tau)|^{2d_\frakg}$.
\item Bismut--Vasserot~\cite{BismutVasserot1989} and
Yoshikawa~\cite{Yoshikawa2004}: $R^{K3}_{\mathrm{Quillen}}$ is a
Borcherds product on the K3 period domain.
\item Conditional scalar witness:
$Z^{1\text{-}\oneloop}_{\BCOV}(K3\times E;\frakg) =
|\eta(\tau)|^{2d_\frakg}|\Phi_{(K3,L)}|^{-2/d_\frakg}V_X^{-d_\frakg/2}$.
\item Oberdieck--Pixton~\cite{OberdieckPixton2018} supplies the
external holomorphic DT datum $\Delta_5(Z)^{-2}$. The one-loop torsion
is compatible with that datum only after additional hCS--Hall/Borcherds
comparison data are supplied.
\end{itemize}

\paragraph{Three attack--heal cycles.}
\textit{Cycle 1.} K\"unneth split is naive at the level of
zeta-determinants because $R_{\mathrm{Quillen}}$ is non-trivial; healed
by Bismut--Gillet--Soul\'e arithmetic-Riemann--Roch.
\textit{Cycle 2.} Ray--Singer Eq.~(7.7) versus Eq.~(8.10) gives a
$(2\pi)^{d_\frakg}$ ambiguity; healed by absorbing into BV-measure
Jacobian and fixing the Quillen gauge to deliver the canonical
$|\eta(\tau)|^{2d_\frakg}$.
\textit{Cycle 3.} Bismut--Vasserot delivers a curvature class only;
healed by Yoshikawa's explicit Borcherds-product evaluation, which
gives the closed K3 dependence. The $\eta(\tau)$ dependence remains an
elliptic Ray--Singer factor unless an explicit determinant-line formula
absorbs it.

\paragraph{What is open.} (i) Full $\tilde q$-expansion match against
$\Delta_5^{-2}$ requires higher Costello--Li loops and compact
hCS--Hall/Borcherds comparison; (ii) the holomorphic trivialisation
relating scalar Quillen metrics to Siegel modular forms is not
constructed here; (iii) on more general CY$_3$ with
$\chi_{\mathrm{top}} \neq 0$, the holomorphic anomaly
$\alpha_{\BCOV} \neq 0$ produces an explicit gauge-anomalous one-loop
counterterm and a deeper Bismut--Quillen anomaly to track.

\section*{References}
\begin{thebibliography}{99}

\bibitem{BCOV1993} M. Bershadsky, S. Cecotti, H. Ooguri, C. Vafa,
\textit{Holomorphic anomalies in topological field theories,}
Nucl.\ Phys.\ B \textbf{405} (1993), 279--304.

\bibitem{BCOV1994} M. Bershadsky, S. Cecotti, H. Ooguri, C. Vafa,
\textit{Kodaira--Spencer theory of gravity and exact results for
quantum string amplitudes,}
Comm.\ Math.\ Phys.\ \textbf{165} (1994), 311--427.

\bibitem{Bismut1986} J.-M. Bismut, \textit{The Atiyah--Singer index
theorems: a probabilistic approach,} J.~Funct.\ Anal.\ \textbf{57}
(1984), 56--99 (and the holomorphic Ray--Singer companion,
Ann.\ Sci.\ ENS \textbf{19} (1986), 305--333).

\bibitem{BGS1988} J.-M. Bismut, H. Gillet, C. Soul\'e,
\textit{Analytic torsion and holomorphic determinant bundles I, II, III,}
Comm.\ Math.\ Phys.\ \textbf{115} (1988), 49--78, 79--126, 301--351.

\bibitem{BismutVasserot1989} J.-M. Bismut, E. Vasserot,
\textit{The asymptotics of the Ray--Singer analytic torsion associated
with high powers of a positive line bundle,}
Comm.\ Math.\ Phys.\ \textbf{125} (1989), 355--367; and follow-up
Comm.\ Math.\ Phys.\ \textbf{125} (1990), 437--469.

\bibitem{CL1201.4501} K. Costello, S. Li,
\textit{Quantization of open-closed BCOV theory I,}
arXiv:1201.4501 (2012).

\bibitem{GriffithsHarris1978} P. Griffiths, J. Harris,
\textit{Principles of Algebraic Geometry,} Wiley, 1978.

\bibitem{OberdieckPixton2018} G. Oberdieck, R. Pixton,
\textit{Holomorphic anomaly equations and the Igusa cusp form
conjecture,} Invent.\ Math.\ \textbf{213} (2018), 507--587.

\bibitem{RaySinger1973} D. B. Ray, I. M. Singer,
\textit{Analytic torsion for complex manifolds,}
Ann.\ of Math.\ \textbf{98} (1973), 154--177.

\bibitem{Yoshikawa2004} K.-I. Yoshikawa,
\textit{$K3$ surfaces with involution, equivariant analytic torsion,
and automorphic forms on the moduli space,}
Invent.\ Math.\ \textbf{156} (2004), 53--117 (and follow-ups
arXiv:math/0211188, arXiv:0708.2008).

\end{thebibliography}

\end{document}
